\usepackage{verbatim}
\usepackage[textsize=scriptsize]{todonotes}
\usepackage{tikz-cd}
\usepackage{etex}
\usepackage{hyperref}
\usepackage{cleveref}
\usepackage[T1]{fontenc}

\usepackage[all,cmtip]{xy}
\usepackage{multicol}


\hypersetup{
   colorlinks,
   linkcolor={red},
   citecolor={blue!50!black},
   urlcolor={blue!80!black}
}

\usepackage{chemarr}
\usepackage{amssymb}
\usepackage{amsmath}
\usepackage{comment}
\usepackage{spectralsequences}
%\usepackage[draft]{spectralsequences}
\usepackage{cleveref}
\usepackage{mathtools}
\usepackage{rotating}
\usepackage{wrapfig}
\usepackage{outlines}
\usepackage{graphicx}
\usepackage{xcolor}
\usepackage{bbm}
\usepackage{enumitem}

%----------------------------------------------------------------------%


\theoremstyle{definition}
\newtheorem{nul}{}[section]
\newtheorem{dfn}[nul]{Definition}
\newtheorem{axm}[nul]{Axiom}
\newtheorem{rmk}[nul]{Remark}
\newtheorem{term}[nul]{Terminology}
\newtheorem{cnstr}[nul]{Construction}
\newtheorem{cnv}[nul]{Convention}
\newtheorem{ntn}[nul]{Notation}
\newtheorem{exm}[nul]{Example}
\newtheorem{obs}[nul]{Observation}
\newtheorem{ctrexm}[nul]{Counterexample}
\newtheorem{rec}[nul]{Recollection}
\newtheorem{exr}[nul]{Exercise}
\newtheorem{wrn}[nul]{Warning}
\newtheorem{qst}[nul]{Question}
\newtheorem*{dfn*}{Definition}
\newtheorem*{axm*}{Axiom}
\newtheorem*{ntn*}{Notation}
\newtheorem*{exm*}{Example}
\newtheorem*{exr*}{Exercise}
\newtheorem*{int*}{Intuition}
\newtheorem*{qst*}{Question}
\newtheorem*{rmk*}{Remark}
\newtheorem*{comp1}{Computation 1:}

\theoremstyle{plain}
\newtheorem{sch}[nul]{Scholium}
\newtheorem{claim}[nul]{Claim}
\newtheorem{thm}[nul]{Theorem}
\newtheorem{prop}[nul]{Proposition}
\newtheorem{lem}[nul]{Lemma}
\newtheorem{var}[nul]{Variant}
\newtheorem{sublem}{Lemma}[nul]
\newtheorem{por}[nul]{Porism}
\newtheorem{cnj}[nul]{Conjecture}
\newtheorem{cor}[nul]{Corollary}
\newtheorem{fact}[nul]{Fact}
\newtheorem*{thm*}{Theorem}
\newtheorem*{prop*}{Proposition}
\newtheorem*{cor*}{Corollary}
\newtheorem*{lem*}{Lemma}
\newtheorem*{cnj*}{Conjecture}
\newtheorem{innercustomthm}{Theorem}
\newenvironment{customthm}[1]
  {\renewcommand\theinnercustomthm{#1}\innercustomthm}
  {\endinnercustomthm}

%%% Make widetilde work in spectral sequences.

\let\oldwidetilde\widetilde
\protected\def\widetilde{\oldwidetilde}


%----------------------------------------------------------------------%

\DeclareMathOperator{\Aut}{\mathrm{Aut}}
\DeclareMathOperator{\Tr}{\mathrm{Tr}}
\DeclareMathOperator{\Res}{\mathrm{Res}}
\DeclareMathOperator{\im}{\mathrm{im}}
\DeclareMathOperator*{\colim}{\mathrm{colim}}
\DeclareMathOperator{\Map}{\mathrm{Map}}
\DeclareMathOperator{\cofiber}{\mathrm{cofiber}}
\DeclareMathOperator{\fiber}{\mathrm{fiber}}
\DeclareMathOperator{\gr}{\mathrm{gr}}
\DeclareMathOperator{\fib}{\mathrm{fib}}
\DeclareMathOperator{\sBar}{\mathrm{Bar}}

\DeclareMathOperator{\Hom}{\mathrm{Hom}} 
\DeclareMathOperator{\End}{\mathrm{End}} 
\DeclareMathOperator{\Skel}{\mathrm{Skel}}
\DeclareMathOperator*{\hocolim}{\mathrm{hocolim}}
\DeclareMathOperator*{\holim}{\mathrm{holim}}
\DeclareMathOperator{\smsh}{\wedge}
\DeclareMathOperator{\coker}{\mathrm{coker}}



\DeclareMathOperator{\Ss}{\mathbb{S}}
\DeclareMathOperator{\FF}{\mathbb{F}}
\DeclareMathOperator{\F}{\mathbb{F}}
\DeclareMathOperator{\ZZ}{\mathbb{Z}}
\DeclareMathOperator{\NN}{\mathbb{N}}
\DeclareMathOperator{\A}{\mathcal{A}}
\DeclareMathOperator{\C}{\mathcal{C}}
\DeclareMathOperator{\D}{\mathcal{D}}
\DeclareMathOperator{\CP}{\mathbb{CP}}
\DeclareMathOperator{\E}{\mathbb{E}}
\DeclareMathOperator{\sE}{\mathrm{E}}
\DeclareMathOperator{\mE}{\mathcal{E}}
\DeclareMathOperator{\MO}{\mathrm{MO}}
\DeclareMathOperator{\Of}{O \langle 4n-1\rangle}
%\DeclareMathOperator{\spOf}{\Sigma^{\infty}_+ O \langle 4n-1\rangle}
%\DeclareMathOperator{\sOf}{\Sigma^{\infty} O \langle 4n-1\rangle}
%\DeclareMathOperator{\Of}{\mathrm{O} \langle 4n-1\rangle}
\DeclareMathOperator{\cS}{\mathcal{S}}
\DeclareMathOperator{\bS}{\mathbb{S}}
\DeclareMathOperator{\Z}{\mathbb{Z}}
\DeclareMathOperator{\Ring}{\textbf{Ring}}
\DeclareMathOperator{\Aff}{\textbf{Aff}}
\DeclareMathOperator{\Ran}{\mathrm{Ran}}
\DeclareMathOperator{\Spec}{\text{Spec}}
\DeclareMathOperator{\Poly}{\text{Poly}}
\DeclareMathOperator{\pic}{\mathrm{pic}}
\DeclareMathOperator{\Pic}{\mathrm{Pic}}
\DeclareMathOperator{\Fin}{\mathrm{Fin}}
\DeclareMathOperator{\Ext}{\mathrm{Ext}}
\DeclareMathOperator{\Gra}{\mathrm{Gr}}
\DeclareMathOperator{\nil}{\text{nil}}
\DeclareMathOperator{\ac}{\tilde{c}}
\DeclareMathOperator{\fin}{\text{fin}}
\DeclareMathOperator{\Gr}{\mathbf{Gr}}
\DeclareMathOperator{\Fil}{\mathbf{Fil}}
\DeclareMathOperator{\Fun}{\text{Fun}}
\DeclareMathOperator{\Alg}{\mathrm{Alg}}
\DeclareMathOperator{\Sp}{\mathrm{Sp}}
\DeclareMathOperator{\Spaces}{\mathcal{S}}
\DeclareMathOperator{\J}{\mathcal{J}}
\DeclareMathOperator{\Cofil}{\mathbf{Cofil}}
\DeclareMathOperator{\GL}{\mathrm{GL}}
\DeclareMathOperator{\cofib}{\mathrm{cofib}}
\DeclareMathOperator{\cof}{\mathrm{cof}}
\DeclareMathOperator{\DD}{\mathbb{cofib}}
\DeclareMathOperator{\Sq}{\mathrm{Sq}}

\newcommand{\spOf}{\Sigma^{\infty}_+ \mathrm{O \langle 4n-1\rangle}}
\newcommand{\sOf}{\Sigma^{\infty} \mathrm{O \langle 4n-1\rangle}}
\newcommand{\wt}{\widetilde}
\newcommand{\BOn}{\mathrm{BO \langle n \rangle}}
\newcommand{\BOm}{\mathrm{BO \langle m \rangle}}
\newcommand{\BO}{\mathrm{BO}}
\newcommand{\Om}{\mathrm{O \langle m -1 \rangle}}
\newcommand{\Oo}{\mathrm{O}}
\newcommand{\U}{\mathrm{U}}
\newcommand{\BP}{\mathrm{BP}}
\newcommand{\Syn}{\mathrm{Syn}}
\newcommand{\HZ}{\mathrm{H}\mathbb{Z}}
\newcommand{\HFp}{\mathrm{H}\mathbb{F}_p}
\newcommand{\HFt}{\mathrm{H}\mathbb{F}_2}
\newcommand{\Qoppa}{\kappa}
\newcommand{\qoppa}{\kappa}
\newcommand{\sampi}{\lambda}
\newcommand{\tor}{\mathrm{tor}}
\newcommand{\tf}{\mathrm{tf}}
\newcommand{\Tot}{\mathrm{Tot}}

\def\o{\mathbbm{1}}
\def\CAlg{\mathrm{CAlg}}
\def\W{\mathbb{W}}
\def\G{\mathbb{G}}
\def\BG{\mathrm{B}\mathbb{G}}
\def\Mfg{\mathcal{M}_{\mathrm{fg}}}
\def\op{\mathrm{op}}
\def\perf{\mathrm{perf}}
\def\P{\mathcal{P}}
\def\CC{\mathbb{C}}

%\DeclareSymbolFont{extraitalic}      {U}{zavm}{m}{it}
%\DeclareMathSymbol{\Qoppa}{\mathord}{extraitalic}{161}
%\DeclareMathSymbol{\qoppa}{\mathord}{extraitalic}{161}
%\DeclareMathSymbol{\sampi}{\mathord}{extraitalic}{166}

% \DeclarePairedDelimiter\abs{\lvert}{\rvert}%
% \makeatletter
% \let\oldabs\abs
% \def\abs{\@ifstar{\oldabs}{\oldabs*}}

%%% Code from Hood to fix a positioning error with his spectralsequences package.
%% \usepackage{etoolbox}
%% \makeatletter
%% \patchcmd\endsseqpage{\path (0,-\sseq@yoffset) node}{\node}{}{\failed}
%% \makeatother


\def\clambda{\Ss\hspace{-2.2pt}/\!\lambda}



%----------------------------------------------------------------------%

\usepackage{tikz}
\usetikzlibrary{matrix,arrows,decorations}
\usepackage{tikz-cd}

\usepackage{adjustbox}

\let\oldtocsection=\tocsection
 
\let\oldtocsubsection=\tocsubsection
 
\let\oldtocsubsubsection=\tocsubsubsection
 
\renewcommand{\tocsection}[2]{\hspace{0em}\oldtocsection{#1}{#2}}
\renewcommand{\tocsubsection}[2]{\hspace{1em}\oldtocsubsection{#1}{#2}}
\renewcommand{\tocsubsubsection}[2]{\hspace{2em}\oldtocsubsubsection{#1}{#2}}

%----------------------------------------------------------------------%
%%% Local Variables:
%%% mode: latex
%%% TeX-master: "Main"
%%% eval: (visual-line-mode t)
%%% eval: (auto-fill-mode 0)
%%% End:
